\section{Ventilation-Perfusion (V/Q) Scan}


\subsubsection*{What Is This Test?}
A V/Q scan has two parts: ventilation (inhaling radioactive gas) and perfusion (injected radioactive material). It compares air flow and blood flow in the lungs. A mismatch (ventilation normal, perfusion absent) suggests pulmonary embolism.

\subsubsection*{Alternative Names}
V/Q scan, lung perfusion scan, lung ventilation scan, ventilation-perfusion scintigraphy, pulmonary V/Q


\subsubsection*{Principle}
The V/Q scan uses two tracers: ventilation imaging uses inhaled radioactive gas (Xe-133 or Kr-81m) or aerosol (Tc-99m DTPA); perfusion imaging uses injected Tc-99m macroaggregated albumin (MAA) that lodges in pulmonary capillaries. A gamma camera images the distribution. Perfusion defects with preserved ventilation indicate pulmonary embolism (mismatch).

\subsubsection*{History}
The V/Q scan was developed in the 1960s using radioactive gases and macroaggregated albumin. It became the standard noninvasive test for pulmonary embolism diagnosis until the widespread adoption of CT angiography.

\subsubsection*{Why This Test Is Ordered}
Your doctor orders a V/Q scan when they suspect a pulmonary embolism (blood clot in your lungs) and a CT scan cannot be performed --- for example, if you have a contrast allergy or kidney disease. It answers the question: is there a mismatch between air flow and blood flow in your lungs that would indicate a blood clot?

\subsubsection*{Is This a Primary or Secondary Test}
A secondary test. CT pulmonary angiography is the preferred first test for suspected pulmonary embolism. V/Q scanning is used when CT is not possible.

\subsubsection*{Interpreting V/Q scan results}
A normal perfusion pattern and a nondiagnostic pattern are not equivalent. The meaning of a result depends on interpretation category, clinical probability, chest findings, and whether SPECT or planar imaging was performed.

\subsubsection*{How This Test Is Performed}
You inhale a small amount of radioactive gas or aerosol through a mask (ventilation part), and receive a small injection of radioactive material into your vein (perfusion part). A special camera then takes images of your lungs. The entire test takes thirty to sixty minutes.

\subsubsection*{What Preparation Is Needed}
No special preparation is needed. Tell your doctor if you are pregnant or breastfeeding.

\subsubsection*{Contraindications and Precautions}
Pregnancy and breastfeeding require examination-specific risk assessment, dose optimization, and current radiopharmaceutical guidance. For FDG studies, follow the protocol for preparation and glucose management.

\subsubsection*{Risks}
Risks associated with nuclear medicine include:
\begin{itemize}
    \item \textbf{Radiation exposure}: Low-dose radiation exposure from the administered radiopharmaceutical tracer.
    \item \textbf{Injection site discomfort}: Minor pain, bruising, or rare infection at the intravenous injection site.
\end{itemize}


You can resume normal activities immediately. Results are typically available within one to three days.

\subsubsection*{Understanding Your Results}

\subsubsection*{Normal Results}
A normal or low-probability result effectively rules out pulmonary embolism.

\begin{center}
\begin{tabular}{ll}
\toprule
\textbf{Population} & \textbf{Reference Range} \\
\midrule
Adult Female & Low probability \\
Adult Male & Low probability \\
Children & Apply age-specific imaging and safety criteria \\
Elderly & Apply age- and comorbidity-specific criteria \\
Pregnancy & Apply pregnancy-specific justification and safety criteria \\
\bottomrule
\end{tabular}
\end{center}
\subsubsection*{Follow-up Tests}
When a test result is not within the normal range, your doctor may recommend additional testing to confirm the finding, identify the underlying cause, and guide treatment decisions:
\begin{itemize}
    \item \textbf{Specific autoantibody panel:} Detailed autoantibody testing based on clinical suspicion (e.g., anti-CCP for rheumatoid arthritis, anti-dsDNA for SLE, ANCA subtypes for vasculitis).
    \item \textbf{Complement levels (C3, C4):} To assess disease activity in lupus (low complement suggests active disease) and evaluate for complement deficiency.
    \item \textbf{Tissue biopsy:} Biopsy of affected organ (skin, kidney, liver, lung) with immunofluorescence staining to confirm autoimmune pathology.
    \item \textbf{Rheumatology consultation:} For comprehensive autoimmune disease evaluation, classification criteria assessment, and treatment planning.
\end{itemize}
\clearpage
